\IfFileExists{stacks-project.cls}{%
\documentclass{stacks-project}
}{%
\documentclass{amsart}
}
\usepackage{amssymb}

% For dealing with references we use the comment environment
\usepackage{verbatim}
\usepackage{fancyvrb}
\usepackage{pmboxdraw}
\newenvironment{reference}{\comment}{\endcomment}
%\newenvironment{reference}{}{}
\newenvironment{slogan}{\comment}{\endcomment}
\newenvironment{history}{\comment}{\endcomment}

% For source figures retained as assets
\usepackage{graphicx}

% For commutative diagrams we use Xy-pic
\usepackage[all]{xy}

% We use 2cell for 2-commutative diagrams.
\xyoption{2cell}
\UseAllTwocells

% We use multicol for the list of chapters between chapters
\usepackage{multicol}

% This is generally recommended for better output
\usepackage{lmodern}
\usepackage[T1]{fontenc}
\usepackage{textalpha}

% For cross-file-references
% We prefer the snapshot tmp/xref/#2.aux,
% which latexmkrc replaces atomically, so
% that documents compiled simultaneously
% never read a half-written .aux file.
% The last argument keeps links to #2.pdf.
\usepackage{xr-hyper}
\newcommand{\maybeexternaldocument}[2]{%
  \IfFileExists{#2.aux}{%
    \externaldocument[#1]{#2}%
  }{%
    \PackageWarning{stack}{Missing #2.aux; external references with prefix #1 unavailable}%
  }%
}
\let\externaldocumentorig\externaldocument
\renewcommand{\externaldocument}[2][]{%
  \IfFileExists{tmp/xref/#2.aux}{%
    \externaldocumentorig[#1]{tmp/xref/#2}[#2.pdf]%
  }{%
  \IfFileExists{#2.aux}{%
    \externaldocumentorig[#1]{#2}%
  }{%
    \PackageWarning{stack}{Missing #2.aux; external references with prefix #1 unavailable}%
  }%
  }%
}

% Package for hypertext links:
\usepackage{hyperref}

% For any local file, say "hello.tex" you want to link to please
% use \externaldocument[hello-]{hello}
\externaldocument[introduction-]{introduction}
\externaldocument[conventions-]{conventions}
\externaldocument[sets-]{sets}
\externaldocument[categories-]{categories}
\externaldocument[topology-]{topology}
\externaldocument[sheaves-]{sheaves}
\externaldocument[sites-]{sites}
\externaldocument[stacks-]{stacks}
\externaldocument[fields-]{fields}
\externaldocument[algebra-]{algebra}
\externaldocument[brauer-]{brauer}
\externaldocument[homology-]{homology}
\externaldocument[derived-]{derived}
\externaldocument[simplicial-]{simplicial}
\externaldocument[more-algebra-]{more-algebra}
\externaldocument[smoothing-]{smoothing}
\externaldocument[modules-]{modules}
\externaldocument[sites-modules-]{sites-modules}
\externaldocument[injectives-]{injectives}
\externaldocument[cohomology-]{cohomology}
\externaldocument[sites-cohomology-]{sites-cohomology}
\externaldocument[dga-]{dga}
\externaldocument[dpa-]{dpa}
\externaldocument[sdga-]{sdga}
\externaldocument[hypercovering-]{hypercovering}
\externaldocument[schemes-]{schemes}
\externaldocument[constructions-]{constructions}
\externaldocument[properties-]{properties}
\externaldocument[morphisms-]{morphisms}
\externaldocument[coherent-]{coherent}
\externaldocument[divisors-]{divisors}
\externaldocument[limits-]{limits}
\externaldocument[varieties-]{varieties}
\externaldocument[topologies-]{topologies}
\externaldocument[descent-]{descent}
\externaldocument[perfect-]{perfect}
\externaldocument[more-morphisms-]{more-morphisms}
\externaldocument[flat-]{flat}
\externaldocument[groupoids-]{groupoids}
\externaldocument[more-groupoids-]{more-groupoids}
\externaldocument[etale-]{etale}
\externaldocument[chow-]{chow}
\externaldocument[intersection-]{intersection}
\externaldocument[pic-]{pic}
\externaldocument[weil-]{weil}
\externaldocument[adequate-]{adequate}
\externaldocument[dualizing-]{dualizing}
\externaldocument[duality-]{duality}
\externaldocument[discriminant-]{discriminant}
\externaldocument[derham-]{derham}
\externaldocument[local-cohomology-]{local-cohomology}
\externaldocument[algebraization-]{algebraization}
\externaldocument[curves-]{curves}
\externaldocument[resolve-]{resolve}
\externaldocument[models-]{models}
\externaldocument[functors-]{functors}
\externaldocument[equiv-]{equiv}
\externaldocument[pione-]{pione}
\externaldocument[etale-cohomology-]{etale-cohomology}
\externaldocument[proetale-]{proetale}
\externaldocument[relative-cycles-]{relative-cycles}
\externaldocument[more-etale-]{more-etale}
\externaldocument[trace-]{trace}
\externaldocument[crystalline-]{crystalline}
\externaldocument[spaces-]{spaces}
\externaldocument[spaces-properties-]{spaces-properties}
\externaldocument[spaces-morphisms-]{spaces-morphisms}
\externaldocument[decent-spaces-]{decent-spaces}
\externaldocument[spaces-cohomology-]{spaces-cohomology}
\externaldocument[spaces-limits-]{spaces-limits}
\externaldocument[spaces-divisors-]{spaces-divisors}
\externaldocument[spaces-over-fields-]{spaces-over-fields}
\externaldocument[spaces-topologies-]{spaces-topologies}
\externaldocument[spaces-descent-]{spaces-descent}
\externaldocument[spaces-perfect-]{spaces-perfect}
\externaldocument[spaces-more-morphisms-]{spaces-more-morphisms}
\externaldocument[spaces-flat-]{spaces-flat}
\externaldocument[spaces-groupoids-]{spaces-groupoids}
\externaldocument[spaces-more-groupoids-]{spaces-more-groupoids}
\externaldocument[bootstrap-]{bootstrap}
\externaldocument[spaces-pushouts-]{spaces-pushouts}
\externaldocument[spaces-chow-]{spaces-chow}
\externaldocument[groupoids-quotients-]{groupoids-quotients}
\externaldocument[spaces-more-cohomology-]{spaces-more-cohomology}
\externaldocument[spaces-simplicial-]{spaces-simplicial}
\externaldocument[spaces-duality-]{spaces-duality}
\externaldocument[formal-spaces-]{formal-spaces}
\externaldocument[restricted-]{restricted}
\externaldocument[spaces-resolve-]{spaces-resolve}
\externaldocument[formal-defos-]{formal-defos}
\externaldocument[defos-]{defos}
\externaldocument[cotangent-]{cotangent}
\externaldocument[examples-defos-]{examples-defos}
\externaldocument[algebraic-]{algebraic}
\externaldocument[examples-stacks-]{examples-stacks}
\externaldocument[stacks-sheaves-]{stacks-sheaves}
\externaldocument[criteria-]{criteria}
\externaldocument[artin-]{artin}
\externaldocument[quot-]{quot}
\externaldocument[stacks-properties-]{stacks-properties}
\externaldocument[stacks-morphisms-]{stacks-morphisms}
\externaldocument[stacks-limits-]{stacks-limits}
\externaldocument[stacks-cohomology-]{stacks-cohomology}
\externaldocument[stacks-perfect-]{stacks-perfect}
\externaldocument[stacks-introduction-]{stacks-introduction}
\externaldocument[stacks-more-morphisms-]{stacks-more-morphisms}
\externaldocument[stacks-geometry-]{stacks-geometry}
\externaldocument[moduli-]{moduli}
\externaldocument[moduli-curves-]{moduli-curves}
\externaldocument[examples-]{examples}
\externaldocument[exercises-]{exercises}
\externaldocument[guide-]{guide}
\externaldocument[desirables-]{desirables}
\externaldocument[coding-]{coding}
\externaldocument[obsolete-]{obsolete}
\externaldocument[fdl-]{fdl}
\externaldocument[index-]{index}

% Heap Project documents
\externaldocument[k3-]{k3}
\externaldocument[torelli-]{torelli}
\externaldocument[linear-algebra-]{linear-algebra}
\externaldocument[groups-]{groups}
\externaldocument[complex-analysis-]{complex-analysis}
\externaldocument[functional-analysis-]{functional-analysis}
\externaldocument[categories-]{categories}
\externaldocument[commutative-algebra-]{commutative-algebra}
\externaldocument[ringed-spaces-]{ringed-spaces}
\externaldocument[homological-algebra-]{homological-algebra}
\externaldocument[lie-algebras-]{lie-algebras}
\externaldocument[representation-theory-]{representation-theory}
\externaldocument[smooth-manifolds-]{smooth-manifolds}
\externaldocument[differential-topology-]{differential-topology}
\externaldocument[algebraic-topology-]{algebraic-topology}
\externaldocument[differential-geometry-]{differential-geometry}
\externaldocument[algebraic-geometry-]{algebraic-geometry}
\externaldocument[symplectic-geometry-]{symplectic-geometry}
\externaldocument[complex-geometry-]{complex-geometry}
\externaldocument[hodge-theory-]{hodge-theory}
\externaldocument[%
variations-of-hodge-structures-]{%
variations-of-hodge-structures}
\externaldocument[%
homological-mirror-symmetry-]{%
homological-mirror-symmetry}
\maybeexternaldocument{deformations-}{deformations}
\externaldocument[vertex-algebras-]{vertex-algebras}
\externaldocument[springer-]{springer}
\externaldocument[infty-categories-]{infty-categories}
\externaldocument[seminars-differential-geometry-]{%
seminars-differential-geometry}
\externaldocument[seminars-complex-geometry-]{%
seminars-complex-geometry}
\externaldocument[seminars-algebraic-geometry-]{%
seminars-algebraic-geometry}
\externaldocument[seminars-analysis-pde-]{%
seminars-analysis-pde}
\externaldocument[surfaces-]{surfaces}
\externaldocument[birational-maps-commutative-algebra-]{birational-maps-commutative-algebra}
\externaldocument[cremona-transformations-]{cremona-transformations}
\externaldocument[derived-categories-of-sheaves-]{derived-categories-of-sheaves}
\externaldocument[geometric-invariant-theory-]{geometric-invariant-theory}
\externaldocument[geometric-stability-conditions-and-group-actions-]{geometric-stability-conditions-and-group-actions}
\externaldocument[moduli-spaces-of-sheaves-]{moduli-spaces-of-sheaves}
\externaldocument[bridgeland-stability-general-theory-]{bridgeland-stability-general-theory}
\externaldocument[hilbert-schemes-]{hilbert-schemes}
\externaldocument[noncommutative-abelian-surfaces-and-kummer-type-hyperkahler-varieties-]{noncommutative-abelian-surfaces-and-kummer-type-hyperkahler-varieties}
\externaldocument[mumford-tate-groups-in-hodge-theory-]{mumford-tate-groups-in-hodge-theory}
\externaldocument[hodge-birational-atoms-2026-]{hodge-birational-atoms-2026}
\externaldocument[xxii-egd-2026-]{%
xxii-egd-2026}
\externaldocument[pde-]{pde}
\externaldocument[probability-]{probability}

\externaldocument[linguistics-]{linguistics}

\externaldocument[%
16th-alga-meeting-2026-]{%
16th-alga-meeting-2026}

\externaldocument[%
iv-jornadas-lefschetz-2026-]{%
iv-jornadas-lefschetz-2026}

\externaldocument[reading-the-masters-]{%
reading-the-masters}

\externaldocument[real-analysis-]{%
real-analysis}

\externaldocument[optimal-transport-]{%
optimal-transport}

\externaldocument[history-]{history}

\externaldocument[seiberg-witten-]{seiberg-witten}

\externaldocument[history-philosophy-biology-]{history-philosophy-biology}

\externaldocument[classical-mechanics-]{classical-mechanics}
\externaldocument[electromagnetism-]{electromagnetism}
\externaldocument[waves-and-optics-]{waves-and-optics}
\externaldocument[wif-]{waves-in-fluids-iv}
\externaldocument[thermal-physics-]{thermal-physics}
\externaldocument[quantum-mechanics-]{quantum-mechanics}
\externaldocument[special-relativity-]{special-relativity}
\externaldocument[physics-]{physics}

\externaldocument[chemistry-]{chemistry}

\externaldocument[molecular-cell-biology-]{molecular-cell-biology}

\externaldocument[microbiology-]{microbiology}

\externaldocument[protistology-]{protistology}

\externaldocument[mycology-]{mycology}

\externaldocument[botany-]{botany}

\externaldocument[zoology-]{zoology}

\externaldocument[animal-histology-]{animal-histology}

\externaldocument[animal-physiology-]{animal-physiology}

\externaldocument[systematics-biogeography-]{systematics-biogeography}

\externaldocument[ecology-natural-resources-]{ecology-natural-resources}

\externaldocument[earth-sciences-]{earth-sciences}

\externaldocument[biostatistics-]{biostatistics}

\externaldocument[mathematical-systems-biology-]{mathematical-systems-biology}

\externaldocument[cimpa-el-salvador-2024-]{cimpa-el-salvador-2024}

\externaldocument[anthropology-]{anthropology}

\externaldocument[yang-mills-theory-]{yang-mills-theory}

\externaldocument[seminars-algebra-]{seminars-algebra}

\externaldocument[seminars-physics-]{seminars-physics}

\externaldocument[seminars-computing-]{seminars-computing}

\externaldocument[oktobermat-2026-]{oktobermat-2026}

\externaldocument[saet-2026-]{saet-2026}

\externaldocument[seminars-symplectic-geometry-]{seminars-symplectic-geometry}

\externaldocument[seminars-dynamics-]{seminars-dynamics}

\externaldocument[oktobermat-2024-]{oktobermat-2024}

\externaldocument[quantum-groups-]{quantum-groups}

\externaldocument[legal-2027-]{legal-2027}

\externaldocument[derived-categories-]{derived-categories}

% Theorem environments.
%
\theoremstyle{plain}
\newtheorem{theorem}[subsection]{Theorem}
\newtheorem{proposition}[subsection]{Proposition}
\newtheorem{lemma}[subsection]{Lemma}
\newtheorem{corollary}[subsection]{Corollary}

\theoremstyle{definition}
\newtheorem{definition}[subsection]{Definition}
\newtheorem{example}[subsection]{Example}
\newtheorem{exercise}[subsection]{Exercise}
\newtheorem{situation}[subsection]{Situation}

\theoremstyle{remark}
\newtheorem{remark}[subsection]{Remark}
\newtheorem{remarks}[subsection]{Remarks}

\numberwithin{equation}{subsection}

% Macros
%
\def\lim{\mathop{\mathrm{lim}}\nolimits}
\def\colim{\mathop{\mathrm{colim}}\nolimits}
\def\Spec{\mathop{\mathrm{Spec}}}
\def\Hom{\mathop{\mathrm{Hom}}\nolimits}
\def\Ext{\mathop{\mathrm{Ext}}\nolimits}
\def\SheafHom{\mathop{\mathcal{H}\!\mathit{om}}\nolimits}
\def\SheafExt{\mathop{\mathcal{E}\!\mathit{xt}}\nolimits}
\def\Sch{\mathit{Sch}}
\def\Mor{\mathop{\mathrm{Mor}}\nolimits}
\def\Ob{\mathop{\mathrm{Ob}}\nolimits}
\def\Sh{\mathop{\mathit{Sh}}\nolimits}
\def\NL{\mathop{N\!L}\nolimits}
\def\CH{\mathop{\mathrm{CH}}\nolimits}
\def\proetale{{pro\text{-}\acute{e}tale}}
\def\etale{{\acute{e}tale}}
\def\QCoh{\mathit{QCoh}}
\def\Ker{\mathop{\mathrm{Ker}}}
\def\Im{\mathop{\mathrm{Im}}}
\def\Coker{\mathop{\mathrm{Coker}}}
\def\Coim{\mathop{\mathrm{Coim}}}

% Boxtimes
%
\DeclareMathSymbol{\boxtimes}{\mathbin}{AMSa}{"02}

%
% Macros for moduli stacks/spaces
%
\def\QCohstack{\mathcal{QC}\!\mathit{oh}}
\def\Cohstack{\mathcal{C}\!\mathit{oh}}
\def\Spacesstack{\mathcal{S}\!\mathit{paces}}
\def\Quotfunctor{\mathrm{Quot}}
\def\Hilbfunctor{\mathrm{Hilb}}
\def\Curvesstack{\mathcal{C}\!\mathit{urves}}
\def\Polarizedstack{\mathcal{P}\!\mathit{olarized}}
\def\Complexesstack{\mathcal{C}\!\mathit{omplexes}}
% \Pic is the operator that assigns to X its picard group, usage \Pic(X)
% \Picardstack_{X/B} denotes the Picard stack of X over B
% \Picardfunctor_{X/B} denotes the Picard functor of X over B
\def\Pic{\mathop{\mathrm{Pic}}\nolimits}
\def\Picardstack{\mathcal{P}\!\mathit{ic}}
\def\Picardfunctor{\mathrm{Pic}}
\def\Deformationcategory{\mathcal{D}\!\mathit{ef}}


\begin{document}

\title{Law and government}
\maketitle
\tableofcontents

\section{The questions of this notebook}
\label{section-gsf-goal}

How is a state organised, what is a right,
and which institutions make and apply the
rules? These are questions of law
(\emph{derecho} in Spanish, \emph{direito} in
Portuguese, \emph{droit} in French), and more
precisely of constitutional law, the part of
law that deals with the organisation of the
state and with fundamental rights.

The notebook first fixes a vocabulary: state
and government, the three public powers,
sources of law, rights and the institutions
that protect them. It then compares the
constitutions of Brazil, Mexico, France, the
United States, China, Russia and Poland
through the same questions. The last part
follows one public policy in detail, the
funding of scientific research: how a
political decision about a budget becomes a
grant, a fellowship or a position.

Constitutional descriptions concern formal
arrangements; they do not establish how
independently institutions act in political
practice. Programme rules and amounts are
dated examples, checked against the primary
sources cited on 26 November 2026.

\section{State, government and public powers}
\label{section-gsf-state}

A useful working distinction is between the
\emph{state}, the continuing public
institutions of a
territory, and the \emph{government}, the
people and
organisation directing public policy at a
particular time.
Changing a cabinet does not recreate the
courts,
universities or public administration.
Usage varies between languages: an English
reference to
``government'' can include institutions that
a Portuguese
discussion would call the \emph{Estado}.
Always ask which meaning a speaker intends.

Three functions provide a first vocabulary.
The \emph{legislative} function makes laws
and authorises
public expenditure. The \emph{executive}
function directs
administration and implements public policy.
The \emph{judicial} function resolves legal
disputes and,
under the relevant system, may review the
legality or
constitutionality of public acts.
These are functions, not a rule that every
country has
three equally powerful and separately elected
bodies.
The constitutional examples below show
different ways
to distribute them.

In Portuguese the terms are \emph{Poder
Legislativo},
\emph{Poder Executivo} and \emph{Poder
Judiciário};
in Spanish, \emph{poder legislativo},
\emph{poder ejecutivo} and \emph{poder
judicial}.
For science policy, the helpful question is
concrete:
who can propose a programme, authorise its
budget,
select its beneficiaries and review a
disputed decision?

\textbf{Head of state and head of
government.}
The first role represents the continuity and
authority
of the state; the second directs the
government.
One person can hold both roles, as in
Brazil's presidency.
Other constitutions distinguish a president
from a prime
minister, but distribute their powers
differently.
The title ``president'' alone therefore tells
us little
about who directs scientific policy.

\section{Law, constitutions and rights}
\label{section-gsf-law-rights}

\textbf{Sources of law.}
Legal rules come from several sources, which
are ordered. At the top is the
\emph{constitution}, which organises the
public powers and limits them. Below it are
the \emph{statutes} (laws) passed by the
legislature, and below these the
\emph{regulations} that the executive issues
in order to apply the statutes. Treaties
between states are a further source, whose
rank depends on the country. A rule of lower
rank that contradicts a higher one can be set
aside; who may do so, and by which procedure,
is one of the main differences between legal
systems.

Two traditions are usually distinguished. In
the \emph{civil law} tradition, to which
France, Brazil and Mexico belong, the main
source is written legislation, organised in
codes. In the \emph{common law} tradition of
England and the United States, earlier
decisions of the courts are themselves
binding in later cases.

\textbf{What is a right?}
A right is a claim of a person that the legal
order recognises and protects. Three things
go together: a holder of the right, someone
(another person or the state) with the
corresponding duty, and a procedure by which
the holder can obtain a remedy when the duty
is not fulfilled. A declaration without the
third element is a political statement more
than a right in the legal sense.

Rights differ by their source. Some come from
a contract, some from a statute, and the
\emph{fundamental rights} from the
constitution itself, so that an ordinary law
cannot remove them. It is usual to
distinguish civil and political rights
(liberty, equality before the law, freedom of
expression, a fair trial, the vote) from
economic, social and cultural rights
(education, health, work, social security).
The Universal Declaration of Human Rights,
adopted by the United Nations General
Assembly in 1948, lists rights of both kinds
\cite{gsf-udhr}. It is a declaration and not
a treaty.

\textbf{Where rights are written.}
Brazil's Constitution lists individual and
collective rights in Article 5 and social
rights in Article 6
\cite{gsf-br-constitution}. Mexico's opens
with human rights and their guarantees, from
Article 1 \cite{gsf-mx-constitution}. In the
United States the first ten amendments, the
Bill of Rights, play this role
\cite{gsf-us-bill-rights}. The preamble of
the French Constitution of 1958 refers to the
Declaration of the Rights of Man and of the
Citizen of 1789 \cite{gsf-fr-constitution}.
The constitutions of Poland, Russia and China
each devote their second chapter to rights
and duties
\cite{gsf-pl-constitution,gsf-ru-constitution,
gsf-cn-constitution}. A written list does not
by itself show how far a right is protected
in practice.

\textbf{Institutions.}
Besides the legislature, the executive and
the ordinary courts, the following kinds of
institution are worth identifying in each
country:
\begin{enumerate}
\item a court or council that can review
whether laws agree with the constitution;
\item the public administration: ministries,
agencies and other bodies that apply the laws
day by day;
\item an audit institution that examines how
public money was spent;
\item an electoral authority;
\item bodies that receive complaints about
violations of rights.
\end{enumerate}
The country sections below name the first
two for each constitution, and the last part
of the notebook names the agencies that fund
research.

\section{Two axes for comparing countries}
\label{section-gsf-comparison}

The first axis concerns territory.
In a \emph{federation}, the constitution
distributes
authority between the federation and
constituent units.
A \emph{unitary state} can still have
substantial local
administration or decentralisation.
The second axis concerns the relationship
between the
executive and legislature: who chooses the
government,
and what can make it leave office?
These axes are independent.

The following is a navigation table, grounded
in the
constitutional articles cited in the country
sections.
It should be read together with their
qualifications.

\begin{center}
\begin{tabular}{p{.20\textwidth}
p{.70\textwidth}}
\textbf{Country} &
\textbf{Formal arrangement to investigate} \\
\hline
Brazil & Federation; president directs the
executive;
Congress has two chambers. \\
Mexico & Federation; president holds
executive power;
Congress has two chambers. \\
France & Unitary, decentralised state;
president and
prime minister; government answerable to the
National Assembly. \\
United States & Federation; president and a
Congress
of two chambers have separate constitutional
roles. \\
China & Unitary state; Communist Party
leadership;
National People's Congress and State Council
within
democratic centralism. \\
Russia & Federation; president directs
overall policy;
government operates under presidential
direction. \\
Poland & Unitary state; president and a
government
whose continuation depends on the Sejm. \\
\end{tabular}
\end{center}

\textbf{Confidence and dismissal.}
In a parliamentary confidence mechanism, a
government
can lose office through a political vote
according to
constitutional rules. That is different from
impeachment
or a procedure for alleged constitutional
misconduct.
In Poland, a constructive vote of no
confidence must
identify a replacement prime minister.
In France, the relevant confidence
relationship is with
the National Assembly, not equally with both
chambers.
Russia's rules allow presidential responses
to a Duma
vote that differ from either example.
The country sections give the precise
references.

\textbf{Formal rules and political practice.}
To study actual power, add evidence about
elections,
party organisation, legislative scrutiny,
judicial
decisions and implementation of budgets.
A constitutional diagram by itself cannot
measure these.
Conversely, a political label cannot tell us
the submission
rules of a research fellowship.

\section{Seven constitutions}
\label{section-gsf-constitutions}

\subsection{Brazil}
\label{section-gsf-brazil-state}

Brazil's Constitution distinguishes
legislative,
executive and judicial powers in Article 2.
The Union, states, Federal District and
municipalities
are autonomous within the constitutional
arrangement
(Articles 1 and 18).
The president exercises executive power with
ministers
(Articles 76 and 84).
Congress consists of the Chamber of Deputies
and Senate
(Articles 44--46); the Supreme Federal Court
has the
constitutional responsibilities in Article
102.
These functions must not be confused with the
internal
governance of a university or funding agency.
See \cite{gsf-br-constitution}.

For budgets, Articles 165--166 connect
executive
initiative with legislative examination.
Distinguish the \emph{Plano Plurianual}
(PPA),
\emph{Lei de Diretrizes Orçamentárias} (LDO)
and \emph{Lei Orçamentária Anual} (LOA):
the multi-year planning framework, budget
guidelines
and annual budget law.
The states have budgets of their own, so a
federal budget figure does not describe all
public spending in the country.

\subsection{Mexico}
\label{section-gsf-mexico}

The Constitution describes a federal republic
(Article 40) and divides federal powers in
Article 49.
Executive power belongs to the president
(Article 80).
Congress consists of the Chamber of Deputies
and Senate
(Article 50), but the annual expenditure
budget is
approved by the Chamber of Deputies under
Article 74 IV.
Constitutional controversies and actions, and
amparo,
have distinct provisions in Articles 103, 105
and 107.
See \cite{gsf-mx-constitution}.

\subsection{France}
\label{section-gsf-france}

France is an indivisible, decentralised
republic
(Article 1). The president's role and
appointment
of the prime minister appear in Articles 5
and 8;
the government's policy role and the prime
minister's
direction of its action appear in Articles
20--21.
Parliament has the National Assembly and
Senate
(Article 24). Government responsibility under
Articles 49--50 concerns the National
Assembly.
The Constitutional Council's review
mechanisms include
Articles 61 and 61-1.
Article 47 regulates parliamentary adoption
of
finance bills, including its special time
limits
\cite{gsf-fr-constitution}.

\subsection{United States}
\label{section-gsf-united-states}

Articles I, II and III of the Constitution
establish
Congress, the presidency and the federal
judiciary.
Congress has the House of Representatives and
Senate;
the executive does not depend on an ordinary
parliamentary confidence vote.
The Tenth Amendment addresses powers reserved
to the
states or people
\cite{gsf-us-constitution,
gsf-us-bill-rights}.
Article I, Section 9 requires appropriations
made
by law for withdrawals from the Treasury.
This is
a key constitutional basis for following
federal
research funding through Congress.
Judicial review is explained through
constitutional
practice, including \emph{Marbury v.
Madison};
it is not simply an explicit sentence in
Article III
\cite{gsf-us-court}.

\subsection{China}
\label{section-gsf-china}

The Constitution states Communist Party
leadership
in Article 1 and democratic centralism in
Article 3.
The National People's Congress is the highest
organ
of state power; it and its Standing Committee
exercise
national legislative power (Articles 57--58).
The Standing Committee is not a second
parliamentary
chamber. The State Council is the highest
organ of
state administration, headed by the premier
(Articles 85--86), and is responsible to the
NPC
and its Standing Committee (Article 92).
Constitutional supervision is assigned within
this
arrangement, including Articles 62 and 67.
It should not be drawn as three equivalent
independent
branches. The NPC approves the state budget
(Article 62); the State Council drafts and
implements it (Article 89)
\cite{gsf-cn-constitution}.

\subsection{Russia}
\label{section-gsf-russia}

The Constitution describes a federation and
separation
of powers (Articles 1, 5 and 10).
The president is head of state and determines
basic
policy directions (Article 80); the
government
exercises executive power under general
presidential
direction (Article 110).
The Federal Assembly consists of the State
Duma and
Federation Council (Articles 94--95).
The president appoints the prime minister
after
Duma approval under Article 111.
Article 117 permits presidential responses to
no-confidence votes; the first such vote does
not
automatically remove the government.
Article 125 concerns the Constitutional
Court.
The government submits the federal budget to
the Duma (Article 114); adopted budget laws
require Federation Council consideration
(Article 106) \cite{gsf-ru-constitution}.

\subsection{Poland}
\label{section-gsf-poland}

Poland's Constitution describes a unitary
state
(Article 3) and distributes legislative,
executive
and judicial powers in Article 10.
The president is head of state (Article 126),
while the Council of Ministers conducts state
policy
under the prime minister's direction
(Articles 146--148).
Parliament consists of Sejm and Senate.
Confidence and government formation involve
the Sejm;
Article 158 requires a no-confidence motion
to name
a replacement prime minister.
Article 188 specifies Constitutional Tribunal
responsibilities \cite{gsf-pl-constitution}.
The Council of Ministers introduces the
budget;
the Sejm adopts it and the Senate can propose
amendments (Articles 219, 221 and 223)
\cite{gsf-pl-budget}.

\section{Science funding}
\label{section-gsf-funding}

Public support for research is a good case
for following a policy from the constitution
to an individual. A legislature authorises
expenditure, a ministry or agency receives
it, and a selection procedure decides which
institution, project or person is supported.

\subsection{From a budget headline to a payment}
\label{section-gsf-budget}

Use the following sequence as a set of
questions,
not as a universal administrative procedure.

\begin{enumerate}
\item \textbf{Proposal.} Who has announced
the figure?
Is it a ministry request, an executive bill
or an
already adopted law?
\item \textbf{Authorisation.} Which
legislature or
other authority approved expenditure, for
which
financial year and which organisation?
\item \textbf{Allocation.} What part reaches
the agency,
university, programme or research project of
interest?
\item \textbf{Selection.} Is there an
institutional
allocation, a grant competition or a
recruitment call?
\item \textbf{Commitment and payment.} Has an
award
actually been made, and when can its
beneficiary
receive the money?
\end{enumerate}

Brazil makes several of these distinctions
explicit
in public-finance law. \emph{Empenho} creates
a spending
obligation subject to its legal conditions;
\emph{liquidação} checks the creditor's right
using the
relevant evidence; an \emph{ordem de
pagamento}
directs payment. These are different stages,
not interchangeable names for a budget
figure.
See Law 4,320/1964, Articles 58 and 62--64
\cite{gsf-br-budget-law}.

\textbf{Does an increase really increase
support?}
First keep the scope and stage constant.
Compare the same institution, programme,
financial year
and price basis. Then separate nominal money
from
purchasing power, and total expenditure from
the part
available for the activity we care about.

For a deliberately fictional example, suppose
a
programme goes from 100 to 110 units while
relevant
prices rise by 10 percent. Its nominal budget
increases,
but the new amount purchases approximately
the same
basket as before. Even without inflation,
more money
for buildings would not imply more fellowship
places.
An unchanged total could support different
numbers of
fellows if the monthly value or funded
duration changes.

This explains why a favourable headline about
the education budget cannot by itself
establish whether the monthly value, the
number or the continuity of a particular
scholarship will change.
Those are three separate questions, each
requiring the
corresponding agency rule, allocation or
award document.

\subsection{Three funding mechanisms}
\label{subsection-gsf-mechanisms}

Institutional support keeps an organisation
operating. A competitive project grant
supports a specified research programme and
can include positions for other researchers.
A personal fellowship supports a named
researcher under its own application and
supervision rules. These mechanisms can
coexist in one department. A research
organisation, which employs researchers,
should also be distinguished from a funding
agency, which distributes money to projects
carried out elsewhere.

\subsection{Brazil: organisations}
\label{section-gsf-brazil-map}

The first map should distinguish a ministry,
a fund,
a funding agency, a host institution and a
researcher.
They perform different jobs even when they
participate
in the same award.

\begin{description}
\item[MCTI and CNPq.]
CNPq is linked to the federal science
ministry and
supports research and researcher training.
Its portfolio includes individual fellowships
and
project competitions \cite{gsf-cnpq}.
\item[MEC and CAPES.]
CAPES is linked to the education ministry,
with a
central role in postgraduate education and
its
evaluation and support \cite{gsf-capes}.
An allocation to a postgraduate programme can
reach
a researcher through a selection run by that
programme.
\item[FNDCT and Finep.]
FNDCT is a science and technology development
fund.
Finep performs its executive-secretariat
functions,
including budgetary and financial
administration
under Law 11,540/2007, Articles 7--9
\cite{gsf-fndct}.
\item[State foundations.]
FAPESP and FAPERJ operate relevant São Paulo
and
Rio de Janeiro routes. Their host, supervisor
and
candidate requirements must be read
separately
\cite{gsf-fapesp-pd,gsf-faperj-pdr10}.
\item[Host, project and fellow.]
The host supplies the institutional setting.
A principal investigator directs a funded
project;
a supervisor supports a particular research
plan.
The fellow receives the fellowship under its
terms.
One person or institution can occupy several
roles.
\end{description}

\subsection{Brazil: postdoctoral fellowships}
\label{section-gsf-brazil-routes}

\textbf{CNPq: read the specific call.}
In the rectified Call 49/2024, the candidate
submitted
the proposal through Plataforma Carlos Chagas
(Section 6.1.1), with the required supervisor
and
institutional consent (Section 6.5.3).
The doctorate-date window was specific to
that edition.
This is an example of a junior postdoctoral
(PDJ) call, not a permanent rule
\cite{gsf-cnpq-pdj}.
The Universal 06/2026 competition also
illustrates a
different route: a principal investigator's
project can
include PDJ or PDS support
\cite{gsf-cnpq-universal2026}.

\textbf{CAPES: institutional allocations.}
PIPD distributes support through institutions
and
postgraduate programmes; the local selection
matters.
Portaria 282/2024 requires an institutional
call
(Article 15). Article 21 includes Brazilian
nationality,
a doctorate obtained at most seven years
before
registration and no employment relationship
with the
institution hosting the postgraduate
programme.
See \cite{gsf-capes-pipd}.
The programme page describes an allocation
period from
October 2024 to September 2027
\cite{gsf-capes-pipd-period}.

PROEX supports highly evaluated postgraduate
programmes
through operating resources and master's and
doctoral
fellowships. Some supported academic
activities can
involve postdoctoral researchers, but a
programme's
PROEX status does not establish a personal
postdoctoral
fellowship for an applicant
\cite{gsf-capes-proex}.

\textbf{FAPESP: a project with a São Paulo
supervisor.}
Under the rules effective from 1 June 2025,
PD proposals
are submitted by the supervisor, with
individual and
project-linked routes. Continuous submission
does not
mean automatic approval.
A candidate can apply before completing the
doctorate,
but must prove completion when accepting the
award.
The rules admit foreign candidates and
normally require
the doctorate to be less than seven years
old, subject
to the specified exceptions.
Read Sections 1, 4--7 for the applicable
conditions
\cite{gsf-fapesp-pd}.

\textbf{FAPERJ: an edition of a competition.}
PDR-10 Call 39/2025 illustrates a
supervisor-led
application in Rio de Janeiro.
It required a doctorate already defended at
submission,
within that edition's specified date window.
Section 2.2.3 also required first-author
publications
or specified innovation alternatives.
In fields where authors are listed
alphabetically, as in mathematics, such a
criterion needs an interpretation by the
agency.
The call permitted Brazilian and foreign
candidates
subject to documentation and other conditions
\cite[Sections 2 and 7]{gsf-faperj-pdr10}.

\textbf{Project vacancies.}
A centre can already hold a grant and then
advertise
a fellowship within it. The applicant
competes for
that position; the centre's principal
investigator
has already obtained the wider project
funding.
The recruitment notice has to be read together
with the funder's rules.

\subsection{Brazil: amounts and conditions}
\label{section-gsf-brazil-amounts}

The following monthly values were checked on
26 November 2026. They are reference values for the named
schemes at that date.

\begin{center}
\begin{tabular}{p{.48\textwidth}
p{.37\textwidth}}
\textbf{Scheme} & \textbf{Monthly value} \\
\hline
CNPq PDJ / PDI (Empresarial) & R\$5,200 \\
CNPq PDS & R\$5,500 \\
CAPES postdoctoral fellowship & R\$5,200 \\
FAPESP PD & R\$12,570 \\
FAPERJ PDR-10, Call 39/2025 & R\$6,500 \\
\end{tabular}
\end{center}

PDI concerns research in companies.
Sources are the agencies' value tables and
the named
FAPERJ call
\cite{gsf-cnpq-rates,gsf-capes-rates,
gsf-fapesp-rates,gsf-faperj-pdr10}.
That PDR-10 edition separately provided
R\$1,000 monthly
for research costs. This is not an additional
unrestricted
personal stipend. Likewise, a technical
reserve or
project travel budget should be recorded
separately.

A useful comparison records the monthly
personal amount,
guaranteed funded duration, renewal
conditions,
research allowance, employment or fellowship
status,
permitted outside work and host location.
The highest monthly number does not answer
all these
questions. Nor does ``up to four years''
necessarily mean
that four years have been awarded at the
beginning.
For PDR-10, the cited edition provided an
initial
twelve months with conditional renewals.

\subsection{Mexico}
\label{subsection-gsf-funding-mexico}

SECIHTI is the federal science secretariat
administering the cited 2026 postdoctoral
programme. Older CONACYT or CONAHCYT
documents should be dated rather than
silently treated
as the current call.
\emph{Estancias Posdoctorales por México
2026}
admitted Mexican and foreign doctoral
graduates for
projects at eligible Mexican public host
institutions.
The candidate's project, adviser and
institutional
support are part of the route; the programme
is not
a salaried position in the secretariat
itself.
Its full-time employment restriction at the
host is
also relevant \cite{gsf-secihti-2026}.

\subsection{France and the European Union}
\label{subsection-gsf-funding-france}

Separate an organisation that conducts
research from
an agency that funds projects. CNRS is a
research
organisation; ANR awards competitive research
funding
\cite{gsf-cnrs,gsf-anr}.
A postdoctoral position in a CNRS-associated
laboratory
does not by itself identify who employs the
researcher
or which grant pays the contract.
Ask for both, and distinguish a fixed-term
project
position from a permanent researcher
competition.

The European Union adds a further funding
route.
The 2026 MSCA European Postdoctoral
Fellowships allowed
applicants of any nationality, applying
jointly with
a host. Global Fellowships had different
nationality
or long-term-residence conditions.
Mobility and other eligibility conditions
still matter
\cite{gsf-msca-2026}.

\subsection{United States}
\label{subsection-gsf-funding-united-states}

Distinguish a federal agency's personal
fellowship
from a university or institute appointment.
The NSF Mathematical Sciences Postdoctoral
Research
Fellowships solicitation 23-603 requires US
citizenship,
national status or permanent residence at
application,
and imposes a specific window since the
doctorate
\cite{gsf-nsf-msprf}.
Eligibility for this programme cannot be
inferred
from eligibility to work at an American
university.

Institute for Advanced Study mathematics
memberships
provide another application route, with their
own
selection and funding arrangements
\cite{gsf-ias}.
For any university postdoc, inspect its
actual notice:
employer, source of support, teaching
obligations,
contract length and candidate eligibility.
Do not transfer the restrictions of one NSF
fellowship
to every position in the country.

\subsection{China}
\label{subsection-gsf-funding-china}

Distinguish the National Natural Science
Foundation
of China (NSFC), which funds research, from
the
Chinese Academy of Sciences (CAS), which
includes
research institutes \cite{gsf-cas}.
NSFC's 2026 Research Fund for International
Scientists
illustrates a project route through an
eligible
mainland host. It required a host employment
contract
before funds were released and substantial
annual
presence in China. A successful project grant
is
therefore not itself a substitute for
securing a
postdoctoral appointment
\cite{gsf-nsfc-rfis-2026}.

The current PIFI 2027 announcement concerns
distinguished
and visiting scientists; it should not be
used as
evidence that an older long-term postdoctoral
category
is still available \cite{gsf-pifi-2027}.

\subsection{Russia}
\label{subsection-gsf-funding-russia}

The Ministry of Science and Higher Education
has
institutional financial-support
responsibilities
\cite{gsf-ru-ministry}.
The Russian Science Foundation (RSF)
separately funds
competitive research projects and teams
\cite{gsf-rsf}.

HSE's 2026--2027 postdoctoral recruitment
supplies
a dated example of institution-level
applications
\cite{gsf-hse-postdoc}.

\subsection{Poland}
\label{subsection-gsf-funding-poland}

Distinguish NCN project funding from NAWA
mobility
and internationalisation programmes.
NCN's SONATA BIS 16 rules describe
postdoctoral
positions planned within a principal
investigator's
project and recruited through an open
procedure
\cite{gsf-ncn-sonata16}.

In ULAM NAWA 2026, the applicant was the
Polish
host institution; the scheme provided for
employment
of a visiting researcher for 6--24 months
\cite{gsf-ulam-2026}.

\section{Questions}
\label{section-gsf-study}

\begin{enumerate}
\item Take one headline about an education
budget in Brazil or France. Identify the
year, the institution, the budget stage and
the baseline of the comparison.
\item Explain why federalism and
parliamentarism answer different questions.
Use Brazil, France and Poland as examples.
\item Compare a French confidence vote,
Poland's constructive vote and the procedure
in Russia. Which constitutional article
determines the next step?
\item For China, explain why treating the
Standing Committee of the National People's
Congress as an upper chamber would
misdescribe the constitutional arrangement.
\item Choose a right written in one of the
seven constitutions. Who holds it, who has
the corresponding duty, and before which
institution can it be claimed?
\item Compare a personal fellowship, a
position paid by a project grant and an
institutional allocation. In each case, who
applies and who selects?
\end{enumerate}

\bibliographystyle{alpha}
\bibliography{my}

\end{document}
